\section{Prompts}
\label{section:appendix_prompt}

\subsection{IRAC KG Generation}
\label{section:appendix_prompt_kg}

Figure \ref{figure:prompt_kg_part1} and Figure \ref{figure:prompt_kg_part2} contain the prompt for producing the IRAC KG. Due to its length, we split it to two parts. We created the KGs using U.S. legal case opinions. Although these cases may contain personally identifying information, they are public data and are searchable via various online tools.
\begin{figure*}[htbp]
\begin{tcolorbox}[enhanced,colback=orange!5!white,colframe=orange!75!black,title=Prompt for Generating IRAC KG (Part I)]

\raggedright

You are tasked with building a knowledge graph by extracting entities and relations from a given legal case using the IRAC (Issue, Rule, Application, Conclusion) framework. This knowledge graph will be used for legal reasoning. Follow these instructions carefully to complete the task.\newline

First, carefully read and analyze the following legal case:\newline

<legal\_case>\newline
\{case\_opinion\}\newline
</legal\_case>\newline

Now, follow these steps to construct the knowledge graph:\newline

1. Entity Extraction:\newline
Extract the following types of entities from the legal case:\newline
- Case: The given case itself\newline
- CitedCase: Other cases cited in the given case\newline
- MaterialFact: Factual elements relevant to the case\newline
- LegalIssue: Legal issues presented in the given case\newline
- Conclusion: The court's ruling on a specific LegalIssue\newline
- Rule: Legal rules or principles applied in the case\newline
- Statute: Laws referenced in the case\newline
- Regulation: Regulations mentioned in the case\newline

For each entity, assign a unique ID, determine its type, and provide a succinct text representing the key aspects of the entity.\newline

2. Relation Extraction:\newline
Identify and extract the following relations between the entities:\newline
- CITES: From Case to CitedCase\newline
- REFERENCES: From Case to Statute or Regulation\newline
- ARISES\_FROM: From LegalIssue to MaterialFact\newline
- ADDRESSES: From Rule to LegalIssue\newline
- APPLIED\_TO: From Rule to MaterialFact\newline
- DERIVES\_FROM: From Rule to CitedCase, Statute, or Regulation\newline
- LEADS\_TO: From Rule to Conclusion\newline

For each relation, assign a unique ID, determine its type, and identify the source (from) and target (to) entities.\newline

\end{tcolorbox}
\caption{Prompt for generating IRAC KG (Part I).}
\label{figure:prompt_kg_part1}
\end{figure*}

\begin{figure*}[htbp]
\begin{tcolorbox}[enhanced,colback=orange!5!white,colframe=orange!75!black,title=Prompt for Generating IRAC KG (Part II)]

\raggedright

3. Output Formatting:\newline
Present your output as a JSON-formatted string that adheres to the following pydantic model structure. Ensure that all required fields (id\_, type\_, label\_ for entities; id\_, type\_, from\_, to\_ for relations) are properly filled.\newline

class Entity(BaseModel):\newline
    id\_: str = Field(description="entity ID")\newline
    type\_: str = Field(description="entity type")\newline
    label\_: str = Field(description="a succinct text representing the key aspects of the entity")\newline

class Relation(BaseModel):\newline
    id\_: str = Field(description="relation ID")\newline
    type\_: str = Field(description="relation type")\newline
    from\_: str = Field(description="source of relation")\newline
    to\_: str = Field(description="target of relation")\newline

class KG(BaseModel):\newline
    vertices\_: List[Entity] = Field(description="vertices from KG")\newline
    relations\_: List[Relation] = Field(description="relations from KG")\newline

4. Ensure that your output is comprehensive, capturing all relevant entities and relations from the given legal case. Be thorough in your analysis and precise in your entity and relation extractions.\newline

5. For entity type "CitedCase", only extract the most important cited cases that are relevant for resolving the legal issues.\newline

6. Do not generate any additional text, explanation or analysis.\newline

\end{tcolorbox}
\caption{Prompt for generating IRAC KG (Part II).}
\label{figure:prompt_kg_part2}
\end{figure*}

\subsection{Prompt for Generating the Explanation Section of the SFT Data}
\label{section:appendix_prompt_sft}

Figure \ref{figure:prompt_sft_part1} and Figure \ref{figure:prompt_sft_part2} show the prompt for producing the SFT training data. Due to its length, we split it to two parts.
\begin{figure*}[htbp]
\begin{tcolorbox}[enhanced,colback=orange!5!white,colframe=orange!75!black,title=Prompt for Generating the SFT Training Data (Part I)]

\raggedright

You are tasked with producing Instruction Finetuning (IFT) data for training a Large Language Model (LLM). The goal is to create a dataset that will teach the model how to identify legal issues and applicable rules based on case facts. This task is crucial for developing AI systems capable of assisting in legal analysis.\newline

First, you will be presented with the facts of a specific legal issue in the legal case. Read these carefully:\newline

<case\_facts>\newline
\{material\_facts\}\newline
</case\_facts>\newline

Analyze these facts thoroughly, considering potential legal issues that might arise from the situation described and potential legal rules that might be applicable to the facts and the legal issues.\newline

Now, here is one actual legal issue identified for this case, and the actual legal rules that are applicable to the given facts and the legal issue.\newline

<legal\_issue>\newline
\{legal\_issue\}\newline
</legal\_issue>\newline

<rules>\newline
\{rules\}\newline
</rules>\newline

Your task is to create IFT data that will teach an LLM to derive this issue from the given facts, and also to identify these legal rules that are applicable to the given case facts and the legal issue. To do this, follow these steps:\newline

1. Write a basic instruction for the LLM without including any details.\newline

2. In your response, include the above case facts, the legal issue, and the identified rules exactly as they were provided to you.\newline

3. Instruct the LLM to provide a brief explanation for the legal issue by relating it back to the given facts, and also on why the rules are applicable to the given case facts and the legal issue.\newline

\end{tcolorbox}
\caption{Prompt for generating the SFT training data (Part I).}
\label{figure:prompt_sft_part1}
\end{figure*}

\begin{figure*}[htbp]
\begin{tcolorbox}[enhanced,colback=orange!5!white,colframe=orange!75!black,title=Prompt for Generating the SFT Training Data (Part II)]

\raggedright

4. Format your output using the following XML format, and do not add any newlines or extra spaces between the XML elements:\newline
\begin{lstlisting}[language=XML, basicstyle=\ttfamily\small, breaklines=true]
<sft_data>
    <sft_input>
        <instruction>
            [Your instruction here]
        </instruction>
        <case_facts>
            [The list of case facts: 
            use separate <fact> tags for each case fact]
        </case_facts>
        <output_format>
            Provide your response in pretty print XML. Use top tag 
            legal_analysis, and the following sub-tags: 
            legal_issue: [The legal issue here]; rules: [A list of rules]; 
            explanation: [Brief explanation of the legal issue by relating 
            it to the case facts, and also on why the rules are applicable 
            to the given case facts and the legal issue].
        </output_format>
    </sft_input>
    <sft_output>
        <legal_issue>
            [The legal issue here]
        </legal_issue>
        <rules>
            [List each applicable rule using the tag rule]
        </rules>
        <explanation>
            [Brief explanation of the legal issue by relating it to the 
            case facts, and also on why the rules are applicable 
            to the given case facts and the legal issue]
        </explanation>
    </sft_output>
</sft_data>
\end{lstlisting}  

Now, create the IFT data for the case facts, the legal issue, and the rules provided earlier. Ensure your instruction is clear and comprehensive, and that your output follows the above XML format.\newline

Do not generate any additional text.\newline

\end{tcolorbox}
\caption{Prompt for generating the SFT training data (Part II).}
\label{figure:prompt_sft_part2}
\end{figure*}

\subsection{Prompt for the LLM Judge during DPO Data Generation}
\label{section:appendix_prompt_dpo}

Figure \ref{figure:prompt_dpo} presents the LLM Judge prompt for producing the DPO training data.
\begin{figure*}[htbp]
\begin{tcolorbox}[colback=orange!5!white,colframe=orange!75!black,title=LLM Judge Prompt for Generating the DPO Training Data]

\raggedright

You are a legal analyst tasked with determining whether specific legal rules can be applied to address a given legal issue.\newline

Your Task:\newline
Analyze whether the provided "Rules to Evaluate" are applicable to the stated legal issue, given the case facts and context.\newline

Information Provided:\newline
Case Facts (the relevant factual circumstances of the case):\newline
\{case\_facts\}\newline

Legal Issue (the specific legal question or dispute that needs to be resolved):\newline
\{legal\_issue\}\newline

Known Applicable Rules (a set of legal rules/statutes/precedents that are already established as relevant to this legal issue):\newline
\{chosen\_rules\}\newline

Rules to Evaluate (another set of legal rules that need to be assessed for applicability):\newline
\{rejected\_rules\}\newline

Your Analysis Should Address:\newline
- Relevance: Do the "Rules to Evaluate" address the same legal subject matter or related issues as the stated legal issue?\newline
- Factual Alignment: Are any factual elements required by these rules present in the given case facts?\newline
- Legal Compatibility: Are these rules consistent with or complementary to the "Known Applicable Rules"?\newline
- Jurisdictional Considerations: Are there any jurisdictional or procedural requirements that affect applicability?\newline
- Scope and Limitations: Are there any known boundaries to how these rules might apply?\newline

Provide your conclusion in JSON format. Use Rules as the top-level attribute; underneath it, include the evaluation for each rule in "Rules to Evaluate" with the following three attributes:\newline
- "Rule": the rule that is being assessed; use the rule text exactly as provided\newline
- "Applicability": Whether each rule in the "Rules to Evaluate" set is applicable (Yes/No/Potentially)\newline
- "Reasoning": Your reasoning for each determination\newline
\end{tcolorbox}
\caption{LLM Judge prompt for generating the DPO training data.}
\label{figure:prompt_dpo}
\end{figure*}
